% En esta sección se escriben las metas del proyecto
\snsection{METAS}
\justifying

Las metas de este proyecto buscan orientar el desarrollo de un sistema didáctico modular que facilite el aprendizaje práctico en electrónica de potencia. Se pretende que los módulos sean fáciles de usar, seguros y adaptables, permitiendo a los usuarios experimentar y comprender mejor los conceptos involucrados. En este sentido, se plantean las siguientes metas específicas:

\begin{itemize}[noitemsep]
    \item Desarrollar un sistema didáctico modular funcional, compuesto por un módulo convertidor reductor, un módulo recortador de onda y un módulo inversor de medio puente.
    \item Lograr que cada módulo pueda ser fácilmente conectado, intercambiado y utilizado dentro del sistema.
    \item Garantizar la seguridad, robustez y facilidad de mantenimiento del sistema modular.
    \item Validar el correcto funcionamiento del sistema completo mediante pruebas con usuarios reales en entorno educativo.
    \item Documentar de forma integral el diseño, funcionamiento y resultados del sistema, incluyendo un manual de usuario complementario.
\end{itemize}

\begin{comment}

Escriba aquí los logros concretos a alcanzar mediante la ejecución del proyecto, las metas se encuentran estrechamente vinculadas con los objetivos.

La meta es un evento futuro hacia el cual dirigimos esfuerzos concretos; es el blanco al que queremos apuntar. Es un punto de llegada.

\begin{itemize}
    \item Son más abarcadoras y generales.
    \item Al redactarse se incluyen todos los aspectos o componentes del proyecto.
    \item Proveen una dirección global del proyecto.
    \item Toman más tiempo en completarse.
    \item Deben ser simples y concisas.
\end{itemize}

\end{comment}
